% Auto-generated by scripts/export_downstream_paper_section.py
\subsection{Downstream Validation}

Overall, downstream validation gives a sharper claim boundary than admission accuracy alone. hard\_v2 supports SkillAdmit-selected utility in a tree-aware coding-agent setting (25/25 versus 23/25), but hard\_v3 does not support universal selected superiority: tree-aware selected ties no\_experience at 29/30, and strict selected is not the best condition. The strongest cross-version result is safety-related: forced harmful artifacts achieve 0/85 success with 85 public-pass/hidden-fail cases, demonstrating systematic negative transfer when harmful experience is executed.

\paragraph{Research questions.}
\begin{itemize}
  \item Does admitted experience improve downstream coding-agent success under hidden-verifier evaluation?
  \item Are repository-tree context and explicit preconditions critical variables for transferring admitted experience?
  \item Does executing harmful admitted experience create systematic negative transfer even when public tests pass?
\end{itemize}

\paragraph{Protocol.}
We evaluate SkillAdmit downstream rather than relying only on admission-label accuracy. hard\_v2 contains 25 import/debug tasks and hard\_v3 contains 30 more agent-like repository repair tasks. Both suites use hidden verifiers, so public-test success is insufficient for success.

The comparison includes no\_experience, SkillAdmit-selected context, selected context with explicit preconditions, precondition-only ablations where available, raw memory, all distilled skills, promoted rules, bad dependency rules, and forced bad artifacts. hard\_v2 and hard\_v3 are reported as distinct evaluation boundaries rather than collapsed into one leaderboard.

\paragraph{Results.}
On hard\_v2, the tree-aware setting gives the clearest positive downstream result for SkillAdmit-selected experience: SkillAdmit-selected solves 25/25 tasks, compared with 23/25 for no\_experience. The selected + precondition context condition also solves 25/25. In contrast, strict visible-file hard\_v2 does not support a selected-win story: selected solves 23/25, while no\_experience solves 24/25 and all distilled skills solve 25/25.

hard\_v3 is the fresh generalization boundary and is deliberately less favorable to a simple selected-superiority claim. In the tree-aware hard\_v3 setting, SkillAdmit-selected and no\_experience both solve 29/30; selected + preconditions also solves 29/30. Raw memory and all distilled skills solve 30/30 and 30/30, respectively.

In strict visible-file hard\_v3, SkillAdmit-selected improves over no\_experience by one task (29/30 versus 28/30), but it is not the best condition: precondition-only and raw memory both solve 30/30. This separates the claim that selected experience can help from the stronger and unsupported claim that selected experience is universally optimal.

The strongest stable safety result is the forced harmful-artifact control. Across hard\_v2 and both hard\_v3 settings, forced bad artifacts solve 0/85 tasks and produce 85 public-pass/hidden-fail cases. This is direct evidence that admitting the wrong experience can systematically damage downstream behavior in ways public tests do not catch.

The ordinary bad\_dependency\_rule rows should not be interpreted as proving that bad advice is safe. In hard\_v3, artifact adherence for that condition is only 4/30 in the tree-aware run and 11/30 in the strict visible-file run, so success often reflects the executor ignoring harmful advice. The forced bad-artifact condition is the cleaner negative-transfer control.

\paragraph{Interpretation.}
These results support a conditional downstream-utility claim: SkillAdmit-selected experience can improve downstream repair success under a coding-agent setting, as shown by hard\_v2 tree-aware validation, but the effect does not generalize into universal dominance on hard\_v3.

The repository tree and explicit preconditions are best treated as context variables rather than as a single universally dominant representation. hard\_v2 shows that precondition text alone is less reliable than tree-aware context, while hard\_v3 shows that precondition-only context can be sufficient for a different task family.

Token savings are not a supported paper claim. The current evidence is about success, hidden-verifier safety, and negative transfer; any token-efficiency claim would require a separate controlled cost study.

\paragraph{Claim limits.}
\begin{itemize}
  \item Do not claim SkillAdmit-selected universally dominates no\_experience.
  \item Do not claim SkillAdmit-selected is the universal best downstream context.
  \item Do not claim cross-version selected token savings.
  \item Do not claim precondition-only generally replaces repository-tree context.
  \item Do not claim raw memory is a safe admission policy from hard\_v3 success alone.
  \item Do not claim bad dependency advice is safe when the model often ignored it.
  \item Do not tune hard\_v2 or hard\_v3 failures while still calling them clean evidence.
\end{itemize}

% Suggested table captions:
% tab:downstream-cross-version: Cross-version downstream validation. hard\_v2 provides positive evidence for tree-aware SkillAdmit-selected utility, while hard\_v3 is a fresh boundary showing that selected context does not universally dominate. Public-hidden counts indicate cases that pass public tests but fail hidden verification.
% tab:negative-transfer: Forced harmful-artifact control. Across hard\_v2 and hard\_v3, forced bad artifacts solve 0/85 tasks and create 85 public-pass/hidden-fail failures.
